\documentclass[11pt,letterpaper]{article}
\usepackage[T1]{fontenc}
\usepackage[utf8]{inputenc}
\usepackage{lmodern}
\usepackage[margin=1in,headheight=15pt]{geometry}
\usepackage{amsmath,amssymb,amsthm,mathtools,mathrsfs}
\usepackage{microtype,booktabs,longtable,array,enumitem}
\usepackage[dvipsnames]{xcolor}
\usepackage{fancyhdr,xurl,needspace}
\usepackage[colorlinks=true,linkcolor=MidnightBlue,citecolor=MidnightBlue,urlcolor=MidnightBlue,breaklinks=true]{hyperref}
\usepackage[nameinlink,noabbrev]{cleveref}
\hypersetup{pdftitle={Vector and Tensor Fields over the Surreal Numbers},pdfauthor={Research exposition prepared for the Surreal project},pdfsubject={An n-dimensional theory of admissible calculus, geometry, and multiscale equations}}
\pagestyle{fancy}\fancyhf{}
\fancyhead[L]{\small Surreal vector and tensor fields}
\fancyhead[R]{\small\thepage}
\renewcommand{\headrulewidth}{0.3pt}
\setlength{\parskip}{0.3em}
\setlength{\parindent}{1.2em}
\setlength{\emergencystretch}{3em}
\setlist{leftmargin=2em,itemsep=0.2em,topsep=0.35em}
\clubpenalty=10000\widowpenalty=10000\displaywidowpenalty=10000
\setcounter{tocdepth}{2}\numberwithin{equation}{section}\allowdisplaybreaks[2]
\newtheorem{theorem}{Theorem}[section]
\newtheorem{proposition}[theorem]{Proposition}
\newtheorem{lemma}[theorem]{Lemma}
\newtheorem{corollary}[theorem]{Corollary}
\theoremstyle{definition}
\newtheorem{definition}[theorem]{Definition}
\newtheorem{example}[theorem]{Example}
\theoremstyle{remark}
\newtheorem{remark}[theorem]{Remark}
\crefname{proposition}{proposition}{propositions}
\crefname{lemma}{lemma}{lemmas}
\crefname{corollary}{corollary}{corollaries}
\crefname{definition}{definition}{definitions}
\crefname{example}{example}{examples}
\newcommand{\No}{\mathbf{No}}
\newcommand{\R}{\mathbb R}\newcommand{\C}{\mathbb C}
\newcommand{\Q}{\mathbb Q}\newcommand{\N}{\mathbb N}\newcommand{\Z}{\mathbb Z}
\newcommand{\A}{\mathscr A}\newcommand{\OO}{\mathcal O}\newcommand{\mm}{\mathfrak m}
\newcommand{\supp}{\operatorname{supp}}\newcommand{\st}{\operatorname{st}}
\newcommand{\id}{\operatorname{id}}\newcommand{\tr}{\operatorname{tr}}
\newcommand{\Ric}{\operatorname{Ric}}\newcommand{\Scal}{\operatorname{Scal}}
\newcommand{\grad}{\operatorname{grad}}\newcommand{\Div}{\operatorname{div}}
\newcommand{\vol}{\operatorname{vol}}\newcommand{\GL}{\operatorname{GL}}
\newcommand{\End}{\operatorname{End}}\newcommand{\Hess}{\operatorname{Hess}}
\newcommand{\Lie}{\mathcal L}\newcommand{\HH}{\mathcal H}
\newcommand{\ev}{\operatorname{Ev}}\newcommand{\jet}{\operatorname{J}}
\newcommand{\halo}[1]{#1^{\mathrm h}}
\newcommand{\abs}[1]{\lvert#1\rvert}\newcommand{\norm}[1]{\lVert#1\rVert}
\newcommand{\vH}{v_{\!H}}
\newcommand{\dd}{\,\mathrm d}
\newcommand{\eps}{\varepsilon}
\newcommand{\repoSHA}{d22a5b35d5b3040e870c3cfd6d1c8f7259e094b0}
\newcolumntype{P}[1]{>{\raggedright\arraybackslash}p{#1}}
\begin{document}
\hypersetup{pageanchor=false}
\begin{titlepage}
\centering
\vspace*{0.25in}
{\large\scshape An $n$-dimensional mathematical framework\par}
\vspace{0.5in}
{\Huge\bfseries Vector and Tensor Fields\par over the Surreal Numbers\par}
\vspace{0.3in}
{\Large Admissible calculus, differential geometry,\par and multiscale field equations\par}
\vspace{0.3in}
{\large September 22, 2026\par}
\vspace{0.35in}
\begin{minipage}{0.94\textwidth}
\begin{abstract}
\noindent
Surreal coefficients extend finite-dimensional tensor algebra without changing
its algebraic identities, but they do not determine a calculus on their own.
This article constructs a theory for every finite dimension $n$ by separating
three settings: finite algebra over a set-sized surreal field; smooth real
spaces with Hahn-valued tensor coefficients; and genuine surreal-coordinate
neighborhoods obtained by support-controlled Taylor prolongation and scale
charts. The principal construction is a sheaf of locally Hahn-supported smooth
fields. We prove that its multivariable Taylor prolongation is well-defined,
respects composition, and identifies coefficient differentiation with the
intrinsic ordered-field derivative on finite halos. Analyticity of the real
coefficient functions is not required.

Within these specified calculi we develop tensors, exterior forms, Lie
brackets, admissible metrics, Levi--Civita connections, curvature, Hodge
operators, divergence, integration and variational equations. Support-preserving
Poincar\'e and Stokes theorems give a compact de Rham base-change theorem. Regular
standard-part reduction commutes with metric geometry; a complementary example
shows that a pointwise invertible metric can have no smooth Hahn inverse and can
carry infinite curvature. Finally, positive-support recursions construct
near-real flows and solutions of nonlinear differential equations whenever a
specified real linearization is invertible. All infinite operations carry
explicit support hypotheses. The framework represents exact hierarchies of
scales, not an automatic regularization of singularities or a transfer of all
real analytic existence theorems.
\end{abstract}
\end{minipage}
\vfill
{\small Prepared in relation to the \texttt{VladimirReshetnikov/Surreal} project.\par
A mathematical exposition with proofs and finite symbolic checks;\par
not a claim of a completed Lean formalization.\par}
\end{titlepage}
\hypersetup{pageanchor=true}
\pagenumbering{roman}
\begingroup\small\setlength{\parskip}{0pt}\tableofcontents\endgroup
\clearpage\pagenumbering{arabic}

\section{Scope and the central choices}\label{sec:scope}

An $n$-dimensional vector with surreal components is an element of $\No^n$,
understood class-theoretically. A vector \emph{field} requires more: its base
space, allowed functions, derivative, coordinate transformations, and any
integration operation must be specified. The two natural questions are
\[
 X:U\subseteq\R^n\longrightarrow\No^n,
 \qquad
 X:W\subseteq\No^n\longrightarrow\No^n.
\]
They are different questions. This article treats both, without identifying
the usual topology of the real domain with a topology inherited from all
surreals.

Throughout, $n\ge1$ is an ordinary finite integer. Indices run from $1$ to
$n$, and repeated upper/lower indices are summed. Tensor rank is finite as
well. The default metric is Euclidean; arbitrary nondegenerate signatures are
introduced explicitly. No step depends on three dimensions unless it is
labeled as such.

\subsection{Three complementary models}

\begin{center}\small
\begin{tabular}{P{0.17\textwidth}P{0.28\textwidth}P{0.43\textwidth}}
\toprule
Model & Base and coefficients & What it supplies\\
\midrule
Finite algebra & $K^n$, with $K\subset\No$ a set-sized real closed field &
Tensors, metrics at a point, contractions, determinants, orthogonal algebra;
no infinite analysis is implicit.\\
Smooth Hahn geometry & An ordinary smooth $n$-manifold $M$, with locally
Hahn-supported smooth coefficients & Coordinate-independent differential
calculus, real-chain integration, global topology and scale-by-scale
geometric equations.\\
Intrinsic coordinate geometry & Open subsets of $K^n$, with rational or
support-certified Taylor--Hahn functions & Actual derivatives with respect
to surreal coordinates, finite halos, arbitrary-center scale charts and
algebraic integration.\\
\bottomrule
\end{tabular}
\end{center}

The main bridge is \cref{thm:prolongation}: a smooth Hahn field has a canonical
Taylor--Hahn realization on a nearstandard open subset of $K^n$. The second
model therefore is not merely symbolic bookkeeping. Its differential
identities become genuine intrinsic ordered-field identities on that subset.
It is nevertheless only one specified class of intrinsic functions.

\subsection{Relation to the repository and the literature}

The repository's current documentation distinguishes finite algebra,
set-sized Hahn workspaces, common-domain coefficient rings, strong summability,
fine differentiation, and the Berarducci--Mantova scalar derivation
\cite{Repo,RepoAnalysis,RepoFound}. We retain those distinctions. In particular,
the one-complex-variable coherent calculus is a precedent for our evaluation
construction; the present treatment extends its \emph{method} to several real
variables, smooth coefficients, tensor sheaves, and geometric equations.

Conway normal forms and real closedness are background results
\cite{Gonshor,vdDE,Poonen93}; the support lemma is classical Hahn--Neumann
algebra \cite{Neumann,Poonen26}. Differential calculus over non-real topological
fields already has a broad theory \cite{BGN}. The tensor identities themselves
are classical differential geometry \cite{Lee,Petersen,Warner}. We give direct
arguments for the support-sensitive extensions rather than claiming that
ordinary completeness or compactness transfers to $K$.

Our contribution is a unified construction with explicit sufficient hypotheses
and proofs of its compatibility results. No claim of priority for every
result, resolution of a named open problem, or exhaustive novelty search is
made. Repository observations are pinned in \cref{sec:audit}; no code in that
repository is modified.

\section{Scalars, size, and admissible infinite sums}\label{sec:scalars}

\subsection{A fixed real closed Hahn workspace}

Write a surreal normal form in increasing-exponent notation as
\begin{equation}\label{eq:normal}
 a=\sum_{\gamma\in S}a_\gamma t^\gamma,
 \qquad a_\gamma\in\R,\qquad t^\gamma:=\omega^{-\gamma},
\end{equation}
where $S\subset\No$ is a well-ordered \emph{set}. Here $\omega^{-\gamma}$
denotes Conway's monomial map, not an unannounced identification with the
Gonshor exponential $\exp(-\gamma\log\omega)$. The relevant law is
$t^{\gamma+\eta}=t^\gamma t^\eta$.

Fix a nonzero, set-sized, divisible ordered additive subgroup
$\Gamma\subset\No$, and put
\begin{equation}\label{eq:workspace}
 K=K_\Gamma=\R((t^\Gamma)).
\end{equation}
The normal-form embedding identifies this field with a subfield of $\No$.
The Hahn-field algebraic-closure theorem applied to $\C((t^\Gamma))=K[i]$
implies that $K$ is real closed \cite{Poonen93}. Thus positive square roots,
finite orthogonal algebra and signature are available in $K$.

\begin{proposition}[Localization of set-sized scalar data]\label[proposition]{prop:localize}
Every set-sized collection of surreal numbers is contained in a workspace
of the form \eqref{eq:workspace}.
\end{proposition}
\begin{proof}
Take the union of the normal-form supports of the given numbers. It is a set.
Its rational additive span in $\No$, enlarged by a nonzero exponent if
necessary, is a set-sized divisible ordered subgroup $\Gamma$. Every original
normal form is a Hahn series over $\Gamma$.
\end{proof}

This proposition localizes scalar data, not arbitrary regularity claims about
functions. It does not assert that every set-valued field has smooth Hahn
coefficients. Nor does it assert that a selected workspace is closed under
every additional scalar derivation.

\subsection{Valuation, order, and standard part}

For $a\ne0$ define $v(a)=\min\supp(a)$, and set $v(0)=+\infty$. Then
\[
 v(ab)=v(a)+v(b),\qquad v(a+b)\ge\min(v(a),v(b)).
\]
The order of $K$ is the sign of the first nonzero real coefficient. Define
\[
 \OO_K=\{a:v(a)\ge0\},\qquad \mm_K=\{a:v(a)>0\}.
\]
These are, respectively, the finite elements and the infinitesimals relative
to $\R$. Every $a\in\OO_K$ has a unique decomposition
\[
 a=\st(a)+h,\qquad \st(a)\in\R,\quad h\in\mm_K,
\]
and coefficient extraction at exponent zero gives a ring homomorphism
$\st:\OO_K\to\R$ with kernel $\mm_K$. Infinite elements have no standard
part under this definition.

\begin{definition}[Strong summability]
A family $(a_j)_{j\in J}$ of Hahn series is strongly summable when the union
of its supports is well ordered and every exponent occurs in the support of
only finitely many $a_j$. Its sum is obtained by finite coefficientwise
addition. The same definition applies to coefficients in any real algebra,
or in a real vector space when only addition is used.
\end{definition}

\begin{lemma}[Classical support lemma]\label[lemma]{lem:neumann}
If $S,T$ are well-ordered subsets of an ordered abelian group, then $S+T$ is
well ordered and a fixed group element has only finitely many representations
$s+t$, with $s\in S,t\in T$. If $T\subset\Gamma_{>0}$ is well ordered, the
additive monoid
\[
 \langle T\rangle=\{0\}\cup\{\tau_1+\cdots+\tau_r:r\ge1,\ \tau_i\in T\}
\]
is well ordered. Every exponent has only finitely many word representations
from $T$, counting all lengths. Consequently, $S+\langle T\rangle$ is well
ordered and all the corresponding coefficient contributions are finite.
\end{lemma}

This is the support theorem imported from Hahn--Neumann theory
\cite{Neumann,Poonen26}. Its last assertion follows by first applying the
finite-pair statement to $S$ and $\langle T\rangle$, then the finite-word
statement to the second summand. We will name these supports whenever a
Taylor substitution, inverse, or nonlinear recursion is used.

\begin{corollary}[Geometric inversion]\label[corollary]{cor:geometric}
If $H$ is a scalar or finite square matrix of Hahn series and the union of
its entry supports is strictly positive, then
\[
 (I-H)^{-1}=\sum_{r\ge0}H^r.
\]
The sum is strong, including all matrix-index choices. The result is an exact
Hahn identity, not a claim about convergence of ordinary partial sums.
\end{corollary}
\begin{proof}
There are finitely many matrix entries. Apply \cref{lem:neumann} to their
support union. At any fixed exponent, only finitely many lengths and index
words contribute. The finite identity
$(I-H)\sum_{r=0}^{N}H^r=I-H^{N+1}$ therefore gives the asserted identity at
each exponent once all words contributing there have been included.
\end{proof}

\subsection{Strong sums are not ordinary limits}

For example, $\sum_{m\ge1}t^{m/(m+1)}$ is a Hahn series over $\Q$.
Its tails have valuations below $1$, so its finite partial sums do not
converge to that series in the valuation topology. Likewise, if
$0<\eta\ll\lambda$ in $\Gamma$, meaning $m\eta<\lambda$ for every
positive integer $m$, then $\sum_{m\ge0}t^{m\eta}$ is strongly summable but
its partial sums are not cofinal in valuation precision. A support argument,
not the phrase ``the terms become small,'' is the justification for our sums.

\section{Finite-dimensional vectors and tensors}\label{sec:algebra}

\subsection{Tensor spaces and transformation laws}

Let $V=K^n$, with basis $e_i$ and dual basis $e^i$. A tensor of type $(r,s)$ is
\[
 T\in T^r_s(V):=V^{\otimes r}\otimes_K(V^*)^{\otimes s},
 \qquad \dim_K T^r_s(V)=n^{r+s}.
\]
Equivalently, it is a multilinear map on $r$ covector inputs and $s$ vector
inputs. Its components are
\[
 T=T^{i_1\cdots i_r}{}_{j_1\cdots j_s}
 e_{i_1}\otimes\cdots\otimes e_{i_r}
 \otimes e^{j_1}\otimes\cdots\otimes e^{j_s}.
\]
Contraction, tensor product, symmetrization, antisymmetrization and permutation
of slots use only finite sums and products. They are therefore valid over
$K$ and, for any given finite data, over the full surreal field.

If $e'_a=e_iP^i{}_a$, with $P\in\GL_n(K)$, then
\begin{equation}\label{eq:transform}
 T'^{a_1\cdots a_r}{}_{b_1\cdots b_s}
 =(P^{-1})^{a_1}{}_{i_1}\cdots(P^{-1})^{a_r}{}_{i_r}
 P^{j_1}{}_{b_1}\cdots P^{j_s}{}_{b_s}
 T^{i_1\cdots i_r}{}_{j_1\cdots j_s}.
\end{equation}
Coordinate transformations $y=\Phi(x)$ instead give vector components
$X_y^a=(\partial y^a/\partial x^i)X_x^i$ and covector components
$\alpha_{y,a}=(\partial x^i/\partial y^a)\alpha_{x,i}$. Distinguishing a
basis change from a coordinate map prevents an inverse-Jacobian ambiguity.

\subsection{Metrics and norms at a point}

A metric at a point is a symmetric nondegenerate bilinear form
$g:V\times V\to K$. Sylvester's law of inertia holds over a real closed
field: in a suitable basis,
\[
 g(u,u)=u_1^2+\cdots+u_p^2-u_{p+1}^2-\cdots-u_n^2,
 \qquad p+q=n.
\]
This follows by symmetric Gaussian elimination and rescaling nonzero diagonal
entries by positive square roots; the numbers of positive and negative terms
are intrinsic by the ordinary finite-dimensional intersection argument.

For positive-definite $g$, put $\norm{u}_g=\sqrt{g(u,u)}\in K_{\ge0}$.
For $v\ne0$, positivity of
$g(u-av,u-av)$ at $a=g(u,v)/g(v,v)$ proves
\[
 g(u,v)^2\le g(u,u)g(v,v).
\]
Consequently the triangle inequality and Gram--Schmidt orthonormalization
hold. These are ordered-field inequalities. The norm is $K$-valued, not a
real-valued Banach norm. No completeness theorem is being asserted.

The metric identifies vectors and covectors by
$X^\flat_i=g_{ij}X^j$ and $\alpha^{\sharp i}=g^{ij}\alpha_j$.
A surreal infinitesimal is not nilpotent: for nonzero $\eps$, every finite
power $\eps^m$ is nonzero. Dropping higher powers therefore changes an exact
tensor into a truncation and must be indicated.

\subsection{Which component scales are invariant?}

The coordinate minimum $v_V(X)=\min_i v(X^i)$ depends on a chosen lattice
$\OO_K^n\subset K^n$. If $P\in\GL_n(\OO_K)$, where the inverse also has
entries in $\OO_K$, both $P$ and $P^{-1}$ cannot lower this minimum. Thus
\[
 v_V(PX)=v_V(X).
\]
The same argument applies to the minimum component valuation of any fixed
finite tensor type. Under arbitrary $\GL_n(K)$ changes this is false:
$e'_1=t^\eta e_1$ changes the first component of $e_1$ to $t^{-\eta}$.
Invariant scale claims should refer either to a specified integral lattice or
to scalar contractions such as $g(X,X)$, $\tr(T^2)$ and curvature invariants.

\section{The sheaf of smooth Hahn fields}\label{sec:sheaf}

\subsection{A common-domain algebra}

For an ordinary open $U\subseteq\R^n$ define
\begin{equation}\label{eq:A}
 \A_\Gamma(U)=C^\infty(U)((t^\Gamma))
 =\left\{\sum_{\gamma\in S}f_\gamma t^\gamma:
 S\subset\Gamma\text{ well ordered},\ f_\gamma\in C^\infty(U)\right\}.
\end{equation}
Support means the exponents with coefficient function not identically zero.
Evaluation at any real $x\in U$ is a well-defined map into $K$.
The representation as a function $U\to K$ is injective: equality at every
$x$ forces equality of each real coefficient function.

\begin{proposition}[Differential algebra]\label[proposition]{prop:diffalgebra}
The algebra \eqref{eq:A} is a commutative $K$-algebra. The operators
\[
 D_i\left(\sum_\gamma f_\gamma t^\gamma\right)
   =\sum_\gamma(\partial_i f_\gamma)t^\gamma
\]
are commuting $K$-linear derivations. They preserve strong summability and
never enlarge support. Pullback along an ordinary smooth map is a strongly
linear algebra homomorphism and satisfies the ordinary chain rule.
\end{proposition}
\begin{proof}
Products have support in $S+T$. By \cref{lem:neumann}, each coefficient is a
finite sum of products of real smooth functions. Differentiating that finite
sum gives Leibniz's rule. Mixed derivatives commute coefficientwise, and
restriction, differentiation and real pullback can only delete exponents.
All the other assertions reduce to finite coefficient operations and the
real smooth chain rule.
\end{proof}

This definition does \emph{not} mean that $U\to K$ is continuous when $U$
has its usual topology and $K$ its intrinsic order topology. It is the chosen
smooth coefficient structure. The intrinsic realization is supplied later.

\subsection{Why local support is necessary}

The presheaf $U\mapsto\A_\Gamma(U)$ need not be a sheaf for arbitrary open
covers. On disjoint intervals $U_m$ one can prescribe the constant section
$t^{1/m}$. Each local section is admissible, but the union of the exponents
$\{1/m:m\ge1\}$ is not well ordered. The glued function has no single
representation \eqref{eq:A} on $\bigcup_mU_m$.

Define $\A^{\mathrm{loc}}_\Gamma$ to consist of functions that admit such a
representation near every real base point. This is a sheaf. Uniqueness of
coefficient extraction makes gluing unambiguous. All local algebra and
calculus statements apply to it. On a smooth manifold, use real coordinate
charts and the pullbacks of \cref{prop:diffalgebra}.

On a compact real manifold there is no global support defect. A finite
subcover of local presentations gives a finite union of well-ordered
supports, hence one well-ordered set. Coefficient functions, or coefficient
tensors, agree on overlaps and glue to global smooth objects. The analogous
statement holds for a section restricted to a neighborhood of a compact set
after a finite localization.

\subsection{Not every nonzero field is a unit}\label{subsec:units}

A sufficient local unit condition is
\begin{equation}\label{eq:unit}
 f=t^\lambda(a+h),\qquad a\in C^\infty(U),\quad a(x)\ne0\text{ on }U,
 \quad \supp(h)\subset\Gamma_{>0}.
\end{equation}
Then
\[
 f^{-1}=t^{-\lambda}a^{-1}
       \sum_{m\ge0}(-a^{-1}h)^m\in\A_\Gamma(U).
\]
If $a>0$, divisibility of $\Gamma$ also gives
\[
 \sqrt f=t^{\lambda/2}\sqrt a
   \sum_{m\ge0}\binom{1/2}{m}(h/a)^m.
\]
All these supports are translates of positive generated monoids.

\begin{example}[A pointwise inverse outside the smooth Hahn algebra]
Fix $\eps=t^\eta$, $\eta>0$, and set $f(x)=x^2+\eps^2$ on a real interval
around zero. Every $f(x)$ is positive and invertible in $K$. Nevertheless,
$f^{-1}$ is not in $\A^{\mathrm{loc}}_\Gamma$ at zero.

Indeed, at $x=0$ its coefficient at exponent $-2\eta$ is $1$. At every
nonzero real $x$ the exact expansion is
\[
 \frac1{x^2+\eps^2}=\sum_{m\ge0}(-1)^m x^{-2m-2}\eps^{2m},
\]
so that same coefficient is zero. It could not be a continuous real
coefficient function on a neighborhood of zero. This example will distinguish
pointwise nondegeneracy from an admissible metric inverse.
\end{example}

\section{Intrinsic derivatives and Taylor--Hahn prolongation}\label{sec:prolong}

\subsection{The topology that is actually used}

Give $K$ its \emph{intrinsic} order topology, whose radii lie in $K_{>0}$.
On $K^n$ use the equivalent coordinate maximum norm
$\norm{x}_\infty=\max_i|x^i|$. A derivative of $F:W\subseteq K^n\to K^m$
at $y$ is a $K$-linear map $L$ satisfying
\begin{equation}\label{eq:fineder}
 \forall\epsilon\in K_{>0}\ \exists\rho\in K_{>0}\quad
 0<\norm{k}_\infty<\rho\ \Longrightarrow\
 \norm{F(y+k)-F(y)-Lk}_\infty<\epsilon\norm{k}_\infty,
\end{equation}
with a ball about $y$ contained in $W$. The linear map is unique: test the
difference of two candidates on small multiples of each coordinate vector.
Polynomial and rational functions are differentiable on their denominator
nonzero domains by finite algebraic remainder identities.

The topology inherited by a set-sized $K\subset\No$ from the \emph{full}
surreal fine topology is different: it is discrete. To see the underlying
obstruction, let $S$ be any set of surreals and fix $x$. There is a positive
surreal smaller than every nonzero distance $|s-x|$, $s\in S$, by the
small-cut property. A sufficiently small ball isolates $x$ from $S\setminus
\{x\}$. Every set-sized subset is thus closed and discrete. In particular,
ordinary real-parameter curves into full $\No$ that are continuous in that
fine topology must be constant. These size-sensitive facts are also emphasized
in the repository \cite{RepoFound}. We never give $K$ that inherited topology
when using \eqref{eq:fineder}.

\subsection{Prolongation of arbitrary smooth coefficients}

For a real open $U\subset\R^n$ put
\[
 \halo U=\{a+h:a\in U,\ h\in\mm_K^n\}\subset K^n.
\]
It is intrinsically open. The decomposition $a=\st(y)$, $h=y-a$ is unique.
For $F=\sum_\gamma f_\gamma t^\gamma\in\A_\Gamma(U)$ define
\begin{equation}\label{eq:eval}
 (\ev F)(a+h)=
 \sum_{\gamma\in\supp(F)}\ \sum_{\alpha\in\N^n}
 \frac{\partial^\alpha f_\gamma(a)}{\alpha!}\,t^\gamma h^\alpha.
\end{equation}
This is a Taylor-\emph{jet} construction. It does not assert convergence of a
real Taylor series to a real smooth function on an ordinary neighborhood.

\begin{theorem}[Multivariable smooth prolongation]\label{thm:prolongation}
Formula \eqref{eq:eval} is strongly summable and defines an injective
$K$-algebra homomorphism into functions on $\halo U$. It extends evaluation
on real points. Its values are intrinsically $C^\infty$ in the sense of
iterated \eqref{eq:fineder}, and
\begin{equation}\label{eq:derbridge}
 \partial_{y^i}(\ev F)=\ev(D_iF).
\end{equation}
For every ordinary smooth map $\Phi:U\to V$,
\begin{equation}\label{eq:natural}
 \ev(\Phi^*F)=(\ev F)\circ(\ev\Phi),
\end{equation}
where $\ev\Phi$ is componentwise prolongation. These statements localize to
$\A^{\mathrm{loc}}_\Gamma$.
\end{theorem}
\begin{proof}
Let $S=\supp(F)$, and let $T$ be the union of the nonzero coordinate supports
of $h$. It is a finite union of well-ordered positive sets. Every summand of
\eqref{eq:eval} has support in $S+\langle T\rangle$. A fixed exponent has
only finitely many choices of $\gamma$, word length, monomials from the
coordinates of $h$, and multiindex: this is \cref{lem:neumann}, with the
finite coordinate labels attached to each word. Thus the displayed double
sum is strong. Evaluation at $h=0$ gives the asserted restriction and proves
injectivity by coefficient extraction on all real points.

For each $a$, taking the full Taylor jet
$\jet_a:C^\infty(U)\to\R[[Z_1,\ldots,Z_n]]$ is an algebra homomorphism.
Leibniz's formula proves this degree by degree. Substitution of positive-order
$h$ into these jets is a homomorphism by the same support lemma. This proves
additivity, multiplicativity and compatibility with $K$ constants.

We give the derivative estimate, since it is the important additional step.
Fix $y=a+h$. For $k\in\mm_K^n$, formal Taylor translation and strong
regrouping yield
\begin{equation}\label{eq:remainder}
 \ev F(y+k)=\ev F(y)+\sum_i\ev(D_iF)(y)k^i
             +\sum_{i,j}k^ik^j R_{ij}(h,k).
\end{equation}
The remainders are evaluations of formal real series with Hahn supports in
$S+\langle T_h\cup T_k\rangle$. If $F\ne0$ and $\beta=\min S$, every
$R_{ij}$ has valuation at least $\beta$, independently of the small $k$.
This follows by removing two of the $k$ factors from terms of $k$-degree at
least two; the remaining factors all have nonnegative valuation. In
particular,
\[
 v\left(\ev F(y+k)-\ev F(y)-\sum_i\ev(D_iF)(y)k^i\right)
 \ge\beta+2v(\norm{k}_\infty).
\]
Given $\epsilon>0$, choose $\rho=t^\theta$ with $\theta>0$ and
$\beta+\theta>v(\epsilon)$. For $0<\norm{k}_\infty<\rho$, the
remainder divided by $\norm{k}_\infty$ has valuation strictly greater than
$v(\epsilon)$ and hence absolute value less than $\epsilon$. The ball stays
in the same halo. This proves \eqref{eq:derbridge}. Applying the same argument
to every $D^\alpha F$ proves existence of all iterated derivatives; its
first-order difference estimate proves their continuity.

Finally, the ordinary smooth chain rule in every finite jet degree gives the
formal identity
\[
 \jet_a(f\circ\Phi)=
 (\jet_{\Phi(a)}f)\circ\bigl(\jet_a\Phi-\Phi(a)\bigr).
\]
The inner series have zero constant terms and evaluate to infinitesimals.
The common positive-support certificate justifies every regrouping, proving
\eqref{eq:natural}. Agreement of local representatives proves localization.
\end{proof}

\begin{remark}[Smooth is not analytic]
If $f(x)=e^{-1/x^2}$ for $x\ne0$ and $f(0)=0$, then $\ev f$ vanishes on
the whole halo $\mm_K$ of zero, because its jet there is zero. It still agrees
with $f$ at every nonzero real point. This is a well-defined smooth
prolongation, not an identity theorem for analytic functions.
\end{remark}

\subsection{Nonuniqueness outside the chosen class}

On $\halo\R=\OO_K$, the functions $y\mapsto y$ and $y\mapsto\st(y)$
agree at every real point. The first has derivative $1$; the second is locally
constant on each open halo and has derivative $0$. Thus restriction to real
points plus intrinsic smoothness does not determine an extension. Theorem
\ref{thm:prolongation} chooses the extension by all real Taylor jets.
The kernel of the derivative is constants on a connected \emph{real base}
inside our Hahn class; it is much larger among all intrinsic smooth functions.

\subsection{Surreal-valued coordinate changes}

Let $\Phi_0:U\to V$ be real smooth and let $H\in\A_\Gamma(U)^m$ have
strictly positive support. Write $\Phi=\Phi_0+H$. For
$F=\sum_\sigma f_\sigma t^\sigma\in\A_\Gamma(V)$ define
\begin{equation}\label{eq:nonrealpull}
 \Phi^*F=\sum_{\sigma,\alpha}
 \frac{(\partial^\alpha f_\sigma)\circ\Phi_0}{\alpha!}
 t^\sigma H^\alpha.
\end{equation}
The support is contained in $\supp(F)+\langle\supp(H)\rangle$.
The proof of \cref{thm:prolongation} applies verbatim to show that pullback is
an algebra map, obeys the chain rule, and realizes actual composition on halos.

\begin{proposition}[Lifting coordinate diffeomorphisms]\label[proposition]{prop:diffeo}
If $\Phi_0$ is a real diffeomorphism and $H$ has positive support, then
$\ev\Phi:\halo U\to\halo V$ is a bijection with an inverse of the same
Taylor--Hahn type. The correction to $\Phi_0^{-1}$ has support in
$\langle\supp(H)\rangle\setminus\{0\}$.
\end{proposition}
\begin{proof}
Write $\Psi=\Phi_0^{-1}+Q$. At exponent $\gamma>0$, the equation
$\Phi\circ\Psi=\id$ is a linear equation for $Q_\gamma$ with coefficient
matrix $D\Phi_0\circ\Phi_0^{-1}$. This matrix has a smooth real inverse.
All remaining terms depend on strictly smaller positive exponents. The
support lemma makes their number finite at each exponent. Transfinite
recursion over the generated monoid constructs a unique $Q$. Construct a
left inverse in the same way. A left inverse and a right inverse coincide by
associativity of the already certified compositions, proving both the formal
and realized statements.
\end{proof}

\section{Genuinely surreal-coordinate neighborhoods}\label{sec:charts}

\subsection{Arbitrary centers and scales}

The finite halo is not all of $K^n$. To work near arbitrary $a\in K^n$, choose
nonzero scales $r_1,\ldots,r_n\in K$ and coordinates
\[
 z^i=\frac{x^i-a^i}{r_i},\qquad z\in\mm_K^n.
\]
One may evaluate every element of
\begin{equation}\label{eq:formalhahn}
 \R[[Z_1,\ldots,Z_n]]((t^\Gamma))
\end{equation}
at such a $z$: its outer Hahn support is fixed, while the coefficient at
each exponent is an arbitrary real formal series. The support proof and the
remainder proof of \cref{thm:prolongation} still apply. In these coordinates
\[
 \partial_{x^i}=r_i^{-1}\partial_{z^i}.
\]
These scale factors matter: an infinitesimal spatial scale can turn a modest
coordinate gradient into an infinite physical-coordinate gradient.

Polynomials and rational functions over $K$ are available independently of
this presentation. Near a point where a rational denominator is nonzero,
shrink a $K$-radius so that its relative change has positive valuation, then
use geometric inversion. Such rational germs are covered by sufficiently
small scale charts even when they fail to be smooth Hahn fields over a large
real base, as in \cref{subsec:units}.

\subsection{An important exclusion}

An arbitrary element of $K[[Z]]$ is not promised to evaluate on the whole
infinitesimal halo. For example,
\[
 \sum_{m\ge0}t^{-m^2\eta}Z^m
 \quad\text{at}\quad Z=t^\eta
\]
has proposed exponents $(-m^2+m)\eta$, not a well-ordered support.
The uniform outer support in \eqref{eq:formalhahn} is a sufficient hypothesis
that excludes this example. Other, larger analytic classes are possible,
but they need their own support certificates or domains. A formal inverse
function theorem alone is not a convergence theorem on a specified domain.

An intrinsic $n$-dimensional atlas may now be built from such open subsets of
$K^n$, requiring every transition and its inverse to belong to an explicitly
admitted calculus. Real diffeomorphism prolongations, their positive-support
perturbations, affine scale changes, and rational local coordinate changes
provide concrete examples. Tensor and connection formulas below apply in
any such atlas once its chain rule and closure hypotheses have been checked.
Global integration and cohomology statements will continue to name the real
base or algebraic chains to which they apply.

\section{Vector fields, tensor fields, and exterior calculus}\label{sec:fields}

\subsection{Modules of geometric fields}

On a real coordinate domain write $\A=\A^{\mathrm{loc}}_\Gamma$ and define
\[
 \mathfrak X_\A(U)=\bigoplus_{i=1}^n\A(U)D_i,
 \qquad
 \Omega^1_\A(U)=\bigoplus_{i=1}^n\A(U)\dd x^i.
\]
These glue under the real Jacobian laws. Tensor fields are sections of the
corresponding finite tensor powers. For an ordinary bundle $E\to M$, its
Hahn extension has locally the same construction; globally a section has
locally well-ordered sums of real sections. We do not identify these
geometric fields with every possible algebraic derivation of a huge function
ring without a locality or strong-linearity condition.

A vector field acts by $X(F)=X^iD_iF$. Its Lie bracket is
\begin{equation}\label{eq:bracket}
 [X,Y]^i=X^jD_jY^i-Y^jD_jX^i.
\end{equation}
Because it is the commutator of derivations, the bracket is antisymmetric and
satisfies Jacobi. Coordinate invariance follows from the chain rule and
cancellation of the symmetric second derivatives of a coordinate change.
It is $K$-bilinear, but not $\A$-bilinear:
$[X,fY]=f[X,Y]+X(f)Y$.

\subsection{Forms, exterior derivative, and Cartan calculus}

Define
\[
 \Omega^p_\A=\bigwedge^p_\A\Omega^1_\A,
 \qquad \operatorname{rank}_\A\Omega^p_\A=\binom np.
\]
For increasing multiindices $I$, set
\[
 d\left(\sum_I\alpha_I\dd x^I\right)
   =\sum_{I,j}D_j\alpha_I\dd x^j\wedge\dd x^I.
\]
Alternation and commutation of the $D_i$ prove $d^2=0$. For homogeneous
$\alpha$ of degree $p$,
\[
 d(\alpha\wedge\beta)=d\alpha\wedge\beta+(-1)^p\alpha\wedge d\beta.
\]
Insertion is $(\iota_X\alpha)(X_2,\ldots,X_p)=\alpha(X,X_2,\ldots,X_p)$.
Define the Lie derivative on forms by Cartan's formula
\begin{equation}\label{eq:cartan}
 \Lie_X=d\iota_X+\iota_Xd.
\end{equation}
It follows algebraically that $[\Lie_X,d]=0$,
$[\Lie_X,\iota_Y]=\iota_{[X,Y]}$ and
$[\Lie_X,\Lie_Y]=\Lie_{[X,Y]}$. To verify these identities directly, note
that both sides are graded derivations and agree on functions and coordinate
one-forms, which generate the exterior algebra. No flow theorem is needed
for this definition.

For a tensor $T$ of type $(r,s)$ the component formula is
\begin{align}\label{eq:lietensor}
 (\Lie_XT)^{i_1\cdots i_r}{}_{j_1\cdots j_s}
 &=X^kD_kT^{i_1\cdots i_r}{}_{j_1\cdots j_s}\notag\\
 &\quad-\sum_{a=1}^r(D_kX^{i_a})
 T^{i_1\cdots k\cdots i_r}{}_{j_1\cdots j_s}
 +\sum_{b=1}^s(D_{j_b}X^k)
 T^{i_1\cdots i_r}{}_{j_1\cdots k\cdots j_s}.
\end{align}
Its signs express the contravariant and covariant transformation laws.
Theorem \ref{thm:prolongation} transports all these identities to halo
realizations, and the intrinsic atlas construction transports them to any
admitted coordinate domain.

\section{Admissible metrics and their Levi--Civita connections}\label{sec:metric}

\subsection{Pointwise versus differential nondegeneracy}

A symmetric matrix $g_{ij}\in\A(U)$ may be nondegenerate at every real
point without its inverse belonging to $\A(U)$. For a \emph{metric in this
calculus}, we require $\det g$ to be a unit of the chosen local coefficient
algebra. When a volume form is used, its square root is required in that
algebra as well. This is an admissibility condition, not an extra condition
on finite-dimensional metric algebra at one point.

A particularly useful class is
\begin{equation}\label{eq:regularmetric}
 g=g_0+h,\qquad g_0\text{ a real smooth nondegenerate metric},
 \quad\supp(h)\subset\Gamma_{>0}.
\end{equation}
Finite matrix inversion gives the exact formula
\begin{equation}\label{eq:metricinverse}
 g^{-1}=\sum_{m\ge0}(-g_0^{-1}h)^m g_0^{-1}.
\end{equation}
Its support is in $\langle\supp(h)\rangle$. The signature is that of $g_0$:
locally diagonalize the real metric smoothly, use nonzero real principal
minors there, and observe that infinitesimal changes cannot alter their signs.
The same argument holds at every point of the prolonged halo.

More generally, constant scale factors are allowed. If $B\in\GL_n(K)$ is
constant and $g=B^T(g_0+h)B$, its inverse remains admissible. This includes
strongly anisotropic metrics not having a nondegenerate real standard part.

\subsection{Construction and uniqueness}

An affine connection satisfies
\[
 \nabla_{fX}Y=f\nabla_XY,\qquad
 \nabla_X(fY)=X(f)Y+f\nabla_XY.
\]
Its coefficients are $\nabla_{D_i}D_j=\Gamma^k{}_{ij}D_k$.
The torsion is $T(X,Y)=\nabla_XY-\nabla_YX-[X,Y]$.

\begin{theorem}[Levi--Civita construction in the admissible algebra]
For every admissible metric there is a unique torsion-free, metric-compatible
connection. It is given by
\begin{equation}\label{eq:christoffel}
 \Gamma^k{}_{ij}=\frac12g^{k\ell}
 (D_i g_{j\ell}+D_j g_{i\ell}-D_\ell g_{ij}).
\end{equation}
\end{theorem}
\begin{proof}
The inverse metric and its derivatives lie in the chosen algebra, so the
formula is defined. It is symmetric in $i,j$, hence torsion-free. Multiplying
by $g_{mk}$ and adding the appropriate permutations gives
$D_i g_{jk}=g_{\ell k}\Gamma^\ell{}_{ij}
+g_{j\ell}\Gamma^\ell{}_{ik}$, which is metric compatibility.
Conversely those three permuted equations, together with symmetry in the
lower indices, force \eqref{eq:christoffel}. This proves uniqueness.
Coordinate independence can either be checked with the chain rule or read
from the coordinate-free Koszul identity
\begin{align*}
 2g(\nabla_XY,Z)&=Xg(Y,Z)+Yg(Z,X)-Zg(X,Y)\\
 &\quad-g(X,[Y,Z])+g(Y,[Z,X])+g(Z,[X,Y]).
\end{align*}
Every term is an operation already established in the differential algebra.
\end{proof}

On general tensors the connection is determined by the product rule and
contraction compatibility:
\begin{align}\label{eq:covtensor}
 (\nabla_kT)^{i_1\cdots i_r}{}_{j_1\cdots j_s}
 &=D_kT^{i_1\cdots i_r}{}_{j_1\cdots j_s}\notag\\
 &\quad+\sum_{a=1}^r\Gamma^{i_a}{}_{k\ell}
 T^{i_1\cdots\ell\cdots i_r}{}_{j_1\cdots j_s}
 -\sum_{b=1}^s\Gamma^\ell{}_{kj_b}
 T^{i_1\cdots i_r}{}_{j_1\cdots\ell\cdots j_s}.
\end{align}
Metric compatibility makes raising and lowering indices commute with
$\nabla$. For a function,
$\Hess_g f{}_{ij}=D_iD_jf-\Gamma^k{}_{ij}D_kf$; the Hessian is symmetric
because the connection has no torsion.

\subsection{Parallel transport and geodesic equations}

Along an admitted curve $x(s)$, parallel transport and geodesics satisfy
\[
 \frac{dY^k}{ds}+\Gamma^k{}_{ij}(x(s))\dot x^iY^j=0,
 \qquad
 \ddot x^k+\Gamma^k{}_{ij}(x(s))\dot x^i\dot x^j=0.
\]
These are well-defined differential equations in an admitted calculus, not
yet universal existence assertions. Metric compatibility and the product
rule imply
\[
 \frac{d}{ds}g(Y,Z)=0,\qquad
 \frac{d}{ds}g(\dot x,\dot x)=0
\]
for parallel fields and geodesics, respectively. These quantities are constant
in the coefficientwise calculus on a connected real parameter interval, and
therefore on its Taylor--Hahn prolongation. In a bare intrinsic calculus,
zero derivative alone is not a global constancy theorem; the example in
Section~5.3 explains the distinction. Existence for positive-support
perturbations of real trajectories is proved in \cref{sec:flows}; exact
polynomial solutions may also be available outside that class.

\section{Curvature and geometric identities}\label{sec:curvature}

Fix the convention
\[
 R(X,Y)Z=\nabla_X\nabla_YZ-\nabla_Y\nabla_XZ-\nabla_{[X,Y]}Z,
 \qquad R(D_i,D_j)D_k=R^\ell{}_{kij}D_\ell.
\]
Then
\begin{equation}\label{eq:curvature}
 R^\ell{}_{kij}=D_i\Gamma^\ell{}_{jk}-D_j\Gamma^\ell{}_{ik}
 +\Gamma^\ell{}_{im}\Gamma^m{}_{jk}
 -\Gamma^\ell{}_{jm}\Gamma^m{}_{ik}.
\end{equation}
In particular $R$ is a tensor although $\Gamma$ is not. Define
\[
 \Ric_{kj}=R^i{}_{kij},\qquad \Scal=g^{kj}\Ric_{kj},\qquad
 G_{ij}=\Ric_{ij}-\tfrac12\Scal\,g_{ij}.
\]

\begin{proposition}[Bianchi identities and conservation]
For an admissible Levi--Civita connection, the first and second Bianchi
identities and the contracted identity hold:
\[
 R(X,Y)Z+R(Y,Z)X+R(Z,X)Y=0,
\]
\[
 (\nabla_XR)(Y,Z)+(\nabla_YR)(Z,X)+(\nabla_ZR)(X,Y)=0,
 \qquad \nabla^iG_{ij}=0.
\]
The curvature has the usual metric symmetries, and $\Ric$ is symmetric.
\end{proposition}
\begin{proof}
Insert \eqref{eq:curvature} into the cyclic sum. Derivative terms cancel
because $\Gamma^\ell{}_{ij}=\Gamma^\ell{}_{ji}$, and the product terms
cancel in pairs. For the second identity, write the connection matrix
$\Gamma=\Gamma_i\dd x^i$ and curvature matrix
$\Omega=d\Gamma+\Gamma\wedge\Gamma$. Direct expansion, using $d^2=0$,
gives
\[
 d\Omega+\Gamma\wedge\Omega-\Omega\wedge\Gamma=0.
\]
This is the covariant second Bianchi identity. Differentiating
$g(Z,W)$ twice and commuting derivatives gives antisymmetry in the two
metric-paired slots; combining it with the first identity gives pair
symmetry and symmetry of Ricci. Contracting the second identity twice and
using $\nabla g=0$ yields
$\nabla^i\Ric_{ij}=\tfrac12D_j\Scal$, hence the last equation.
All operations are finite tensor contractions in the admissible algebra.
\end{proof}

In dimension two $G$ vanishes identically. In dimensions at least three it
need not. Flatness means $R=0$, not that the coefficients happen to be small.
A nonzero infinitesimal curvature is still nonzero curvature.

\subsection{First-order metric response with an exact remainder}

Let $g=g_0+t^\eta h$, $\eta>0$, with $h$ a real symmetric tensor. The
coefficient at $t^\eta$ of the connection difference is the tensor
\begin{equation}\label{eq:variationGamma}
 C^k{}_{ij}=\tfrac12g_0^{k\ell}
 (\nabla^0_i h_{j\ell}+\nabla^0_j h_{i\ell}-\nabla^0_\ell h_{ij}).
\end{equation}
Subtracting the metric-compatibility equations for the two connections and
extracting the first coefficient proves this formula. Substitution into
\eqref{eq:curvature} gives
\begin{equation}\label{eq:variationR}
 R^\ell{}_{kij}(g)=R^\ell{}_{kij}(g_0)
 +t^\eta(\nabla^0_i C^\ell{}_{jk}-\nabla^0_j C^\ell{}_{ik})
 +E^\ell{}_{kij},\qquad \vH(E)\ge2\eta.
\end{equation}
Here $\vH$ is the minimum nonzero \emph{coefficient-function} exponent in
a fixed local real frame. The remainder is a genuine Hahn tensor. Formula
\eqref{eq:variationR} is an exact decomposition, rather than an equality in
a ring where $(t^\eta)^2$ is zero.

For $h=\Lie_Vg_0$, naturality under the prolonged diffeomorphisms implies
that the first curvature response is $\Lie_VR(g_0)$. This identifies the
usual infinitesimal coordinate freedom and explains why solving geometric
field equations requires a treatment of gauge, not just componentwise
inversion of a linear operator.

\section{Volume, Hodge duality, and vector calculus in every dimension}\label{sec:hodge}

\subsection{Volume and the Hodge star}

For an oriented admissible metric of signature $(p,q)$ set
\[
 \mu_g=\sqrt{|\det g|},\qquad
 \vol_g=\mu_g\dd x^1\wedge\cdots\wedge\dd x^n.
\]
For \eqref{eq:regularmetric}, the square root is obtained by the positive
binomial expansion after factoring $|\det g_0|$. The induced inner product
on forms is the determinant pairing. Define $*$ by
\begin{equation}\label{eq:star}
 \alpha\wedge *\beta=\langle\alpha,\beta\rangle_g\vol_g,
 \qquad *:\Omega^k_\A\to\Omega^{n-k}_\A.
\end{equation}
An oriented pseudo-orthonormal basis gives
\begin{equation}\label{eq:star2}
 *^2\alpha=(-1)^{k(n-k)+q}\alpha.
\end{equation}
Indeed the two permutations contribute $(-1)^{k(n-k)}$, and the product of
all metric signs is $(-1)^q$. This proves the formula without a positivity
assumption.

\subsection{Gradient, divergence, and Laplacians}

The definitions and their coordinate formulas are
\begin{align}
 \grad_g f&=(df)^\sharp=g^{ij}D_jf\,D_i,\label{eq:grad}\\
 \Lie_X\vol_g&=(\Div_gX)\vol_g,
 &\Div_gX&=\mu_g^{-1}D_i(\mu_gX^i)=\nabla_iX^i,\label{eq:div}\\
 \Delta_gf&=\Div_g\grad_gf
 =\mu_g^{-1}D_i(\mu_g g^{ij}D_jf).\label{eq:lap}
\end{align}
For \eqref{eq:div}, use Cartan's formula and expand
$d\iota_X\vol_g$. Equality with the connection trace follows from
$D_i\det g=(\det g)\tr(g^{-1}D_i g)$ and
$\Gamma^j{}_{ji}=D_i\log\mu_g$, with the latter notation meaning the
algebraic quotient $\mu_g^{-1}D_i\mu_g$, whether or not a logarithm has
been selected in the function ring.

On $k$-forms define the codifferential by
\begin{equation}\label{eq:codiff}
 \delta_g=(-1)^{n(k+1)+q+1}*d*,\qquad
 \Delta_H=d\delta_g+\delta_gd.
\end{equation}
The sign follows by integrating $d(\alpha\wedge*\beta)$ on real compact
chains, or by the equivalent local algebraic identity. On functions
$\Delta_H=-\Delta_g$. This distinguishes the nonnegative real Riemannian
Hodge-Laplacian convention from the positive-coordinate-second-derivative
convention in \eqref{eq:lap}.

\subsection{What replaces curl?}

The dimension-independent vorticity is the two-form
\[
 \mathcal W_X=d(X^\flat).
\]
Its dual $*\mathcal W_X$ has degree $n-2$. It is a scalar in dimension two,
a one-form (hence a vector after raising an index) in dimension three, and a
two-form in dimension four. In dimension one it vanishes. Thus a vector-valued
curl is not a dimension-independent operation.

Similarly, the metric and orientation give an $(n-1)$-ary cross product
\[
 X_1\times\cdots\times X_{n-1}
 =\bigl(*(X_1^\flat\wedge\cdots\wedge X_{n-1}^\flat)\bigr)^\sharp.
\]
It is the familiar binary cross product only when $n=3$. Extra exceptional
algebraic structures in other dimensions are not supplied by merely changing
the scalar field.

The identities $d^2=0$ and $\delta_g^2=0$ replace the usual ``curl of a
gradient'' and ``divergence of a curl'' formulas. In a flat constant metric,
\[
 \delta_gd(X^\flat)=d(\Div_gX)-\Delta_g(X^\flat),
\]
where the final Laplacian is taken componentwise. For a variable metric the
Hodge identity remains valid, while rewriting it using a rough covariant
Laplacian introduces the usual Ricci curvature term. The flat formula must
not be transplanted unchanged to curved coordinates or a curved metric.

\section{Integration, Stokes, and local potentials}\label{sec:integration}

\subsection{Integration on real chains is coefficientwise}

For a smooth real oriented compact $p$-chain $C$ and a Hahn $p$-form defined
near it, use a finite localization and define
\begin{equation}\label{eq:integral}
 \int_C^{H}\sum_\gamma\alpha_\gamma t^\gamma
 =\sum_\gamma\left(\int_C\alpha_\gamma\right)t^\gamma.
\end{equation}
The sum is well defined because its support is a subset of a single
well-ordered support. Real integration happens \emph{inside each coefficient}
first. Formula \eqref{eq:integral} is not defined as a fine limit of
surreal-valued Riemann sums. It is independent of finite real chart choices
and is $K$-linear by the finite-pair support lemma.

\begin{theorem}[Hahn Stokes theorem]\label{thm:stokes}
For $C$ a smooth real compact oriented $p$-chain and
$\alpha\in\Omega^{p-1}_\A$ defined near it,
\[
 \int_C^{H}d\alpha=\int_{\partial C}^{H}\alpha.
\]
\end{theorem}
\begin{proof}
Every coefficient form satisfies real Stokes. Taking those equalities as the
coefficients of one Hahn series proves the identity. No exchange of an
uncertified infinite integral with an infinite sum is involved: the integral
was defined coefficientwise on a common support.
\end{proof}

For a real domain with smooth compact boundary this yields
\[
 \int_\Omega^{H}\Div_gX\,\vol_g
 =\int_{\partial\Omega}^{H}\iota_X\vol_g,
\]
provided the metric, inverse, density and field are admissible near the
closure. The flux form is the primary definition; a unit-normal expression
requires the corresponding induced metric and square roots.

\begin{proposition}[Positivity for smooth common-support integrands]
Let $M$ be a nonempty compact real manifold without boundary, with a positive
real smooth density $\nu$. If $f\in\A_\Gamma(M)$ is pointwise nonnegative
at every real point and is not the zero field, then $\int_M^H f\nu>0$.
\end{proposition}
\begin{proof}
Let $\gamma_0$ be the least exponent whose smooth coefficient function is
not identically zero. Pointwise nonnegativity forces
$f_{\gamma_0}(x)\ge0$ everywhere: wherever it is nonzero it is the leading
coefficient of $f(x)$. It is positive on a nonempty real open set. Hence its
real integral is positive, and this is the leading coefficient of the Hahn
integral. The same argument works for compact domains with interior whenever
the nonzero smooth field is detected there.
\end{proof}

This result concerns smooth common-support densities on real domains. It is
not a statement about arbitrary Hahn-valued measures or a change from
coefficientwise to strong countable additivity.

\subsection{A support-preserving Poincar\'e lemma}

Suppose $U\subset\R^n$ is star-shaped about zero. For a real $p$-form,
$p\ge1$, define the radial homotopy
\[
 (H\alpha)_x(v_1,\ldots,v_{p-1})
 =\int_0^1 s^{p-1}\alpha_{sx}(x,v_1,\ldots,v_{p-1})\dd s.
\]
Extend it coefficientwise to Hahn forms, and put $H=0$ on functions.

\begin{theorem}[Support-preserving local exactness]\label{thm:poincare}
For common-support forms in $\Omega^p(U;\R)((t^\Gamma))$, the homotopy
never enlarges Hahn support and satisfies
\[
 dH+Hd=\id-e_0^*,
\]
where $e_0:U\to U$ is the constant map. Consequently every closed positive-
degree common-support Hahn form on $U$ is exact with a primitive having
no additional exponents. This proves local exactness for the sheaf as well. A Hahn function with all $D_i f=0$ is constant on every connected
real domain.
\end{theorem}
\begin{proof}
For the dilation $r_s(x)=sx$, differentiating $r_s^*\alpha$ in $s$ and
using Cartan's identity gives a total derivative whose integral from $0$ to
$1$ is $\alpha-e_0^*\alpha$. This is the displayed real homotopy identity;
each coefficient is smooth at $s=0$ because the factor $s^{p-1}$ is
nonnegative-degree. Applying the identity coefficientwise proves it for
Hahn forms. For a positive-degree form $e_0^*\alpha=0$. Finally, a real
coefficient function with zero gradient is constant on a connected real
open set, and the same conclusion holds at every Hahn exponent.
\end{proof}

In particular, a curl-free field has a scalar potential locally precisely
when $d(X^\flat)=0$; ``curl-free'' here means the two-form condition in every
dimension. A divergence-free field satisfies
$d(\iota_X\vol_g)=0$ and locally has an $(n-2)$-form potential when $n\ge2$.
These are differential-form statements, not a special three-dimensional
vector-potential ansatz.

\subsection{Integration on surreal-coordinate simplices}

There is a second, algebraic integration theory that does not require a real
base. For polynomial forms over $K$ on an oriented affine $p$-simplex with
vertices in $K^n$, pull back to barycentric coordinates
$\lambda_0+\cdots+\lambda_p=1$ and define its moments by
\begin{equation}\label{eq:simplex}
 \int_{\Delta^p}^{\mathrm{alg}}
 \lambda_0^{a_0}\cdots\lambda_p^{a_p}
 \dd\lambda_1\cdots\dd\lambda_p
 =\frac{a_0!\cdots a_p!}{(p+a_0+\cdots+a_p)!}.
\end{equation}
Extend by $K$-linearity and the oriented affine pullback. Polynomial Stokes
holds. One proof applies the polynomial fundamental theorem successively in
simplex coordinates. Another observes that, at each fixed degree, both sides
are polynomial expressions in the coefficients and vertices with rational
coefficients; their equality for all real data is a polynomial identity and
hence remains true over $K$.

A support-certified formal form on an infinitesimal scaled simplex can also
be integrated termwise by these moments. The same positive-support proof
justifies its sums. This supplies exact small-scale integrals, but neither
construction defines a countably additive Lebesgue theory on every subset of
$K^n$. Classical global integral theorems require the chain category stated
in their hypotheses.

\section{Global cohomology and a precise Hodge theorem}\label{sec:cohomology}

Let $M$ be a compact smooth real manifold without boundary. Denote by
$\Omega_H^p(M)$ the global sections of the Hahn form sheaf. Compactness of
the \emph{real} base gives
\[
 \Omega_H^p(M)=\Omega^p(M;\R)((t^\Gamma)).
\]
This is a statement about supports of real coefficient forms, not compactness
of the much larger halo in its intrinsic topology.

\begin{theorem}[Compact de Rham base change]\label{thm:derham}
There is a canonical isomorphism of graded $K$-algebras
\begin{equation}\label{eq:cohom}
 H^\bullet(\Omega_H^\bullet(M),d)
 \cong H^\bullet_{\mathrm{dR}}(M;\R)\otimes_\R K.
\end{equation}
In particular, the dimension in degree $p$ is the ordinary real Betti number
$b_p(M)$.
\end{theorem}
\begin{proof}
The real cohomology in each degree is finite-dimensional. Choose real closed
forms $\theta_1,\ldots,\theta_b$ representing a basis. Choose real linear
coordinate maps on closed forms and a linear choice of primitive for exact
forms; such a choice exists by a vector-space splitting. Let $P$ apply that
primitive choice after subtracting the selected cohomology representative.
Thus every closed
real form has an identity
\[
 \alpha=\sum_{j=1}^b c_j(\alpha)\theta_j+dP\alpha,
\]
where each $c_j$ and $P$ is real linear. Apply these fixed operators to every
coefficient of a closed Hahn form. They cannot add exponents. We obtain
\[
 \alpha=\sum_{j=1}^b a_j\theta_j+d\beta,
 \qquad a_j\in K,\quad\beta\in\Omega_H^{p-1}(M).
\]
This proves surjectivity of the canonical map from the right side of
\eqref{eq:cohom}. If a finite $K$-linear combination of the $\theta_j$ is
exact, coefficient extraction gives a real exact combination at every
exponent. Real independence forces all its coefficients to vanish, proving
injectivity. The map itself sends scalar-extended real cohomology classes to
the same representative forms, so it is independent of the auxiliary
splittings. Products agree because wedge products have finite coefficient
convolutions.
\end{proof}

For an oriented compact $M$ choose an ordinary real Riemannian metric $g_0$.
The classical real Hodge theorem \cite{Warner} supplies harmonic projection
$P_0$ and Green operator $G_0$ on smooth forms, with
\[
 \id=P_0+d\delta_0G_0+\delta_0dG_0.
\]
Every operator in this identity extends coefficientwise without enlarging
support. We immediately obtain the following exact, but carefully scoped,
Hodge theorem.

\begin{corollary}[Coefficientwise Hodge decomposition]\label[corollary]{cor:hodge}
For the real metric $g_0$,
\[
 \Omega_H^p(M)=d\Omega_H^{p-1}(M)
 \oplus\delta_0\Omega_H^{p+1}(M)
 \oplus\bigl(\HH_{g_0}^p(M;\R)\otimes_\R K\bigr).
\]
The summands are orthogonal for the $K$-valued pairing
$\int_M^H\alpha\wedge *_0\beta$. This pairing is positive definite.
Every Hahn cohomology class has a unique harmonic representative.
\end{corollary}
\begin{proof}
Extend the real operator identities and orthogonality coefficientwise.
Uniqueness follows by applying the real projections at each exponent.
For positivity, the leading coefficient of the pointwise squared norm of a
nonzero form is the real squared norm of its leading nonzero coefficient
form; its integral is positive. Finite dimensionality of the harmonic space
identifies its Hahn extension with tensoring by $K$.
\end{proof}

This does not claim an unrestricted Hodge theorem for every surreal metric,
every intrinsic $K$-manifold, or a completed infinite-dimensional normed
space. The fixed real elliptic theory is the input. Establishing a different
metric's global Green operator requires a separate argument.

For example, on the real $n$-torus,
$\dim_K H_H^p=\binom np$. On a circle, $a\dd\theta$ has period $2\pi a$
and is nonexact for every $a\in K\setminus\{0\}$, even when $a$ is
infinitesimal. Scales enrich the coefficients of topological classes; in this
model they do not create extra Betti numbers. This cohomology is not asserted
to be singular cohomology of the halo with its fine, highly disconnected
topology.

\section{Standard-part geometry and scale stability}\label{sec:reduction}

Let $\A_{\ge0}$ be the subalgebra with no negative exponents, and
$\A_{>0}$ its ideal of positive-support sections. Coefficient extraction
extends standard part to
\[
 \st:\A_{\ge0}\to C^\infty(M),\qquad
 \ker\st=\A_{>0}.
\]
The maps $D_i$ commute with it because they act only on coefficient functions.

\begin{theorem}[Regular reduction of tensor geometry]\label{thm:reduction}
Suppose $g\in\A_{\ge0}$ is symmetric and $g_0=\st(g)$ is a real smooth
nondegenerate metric. For tensor fields with nonnegative support,
standard part commutes with tensor products, contractions, exterior
products, $d$, Lie brackets and Lie derivatives. Moreover,
\[
 \st(g^{-1})=g_0^{-1},\quad
 \st(\Gamma(g))=\Gamma(g_0),\quad
 \st(R(g))=R(g_0),\quad
 \st(\Ric(g))=\Ric(g_0),
\]
and similarly for scalar curvature, volume density, Hodge star,
gradient and divergence. Coefficientwise real-chain integrals commute with
standard part whenever the integrand has nonnegative support.
\end{theorem}
\begin{proof}
All finite polynomial expressions commute with the ring homomorphism $\st$,
and it commutes with $D_i$. Formula \eqref{eq:metricinverse} shows that
$g^{-1}$ has no negative exponents and has constant coefficient $g_0^{-1}$.
The binomial density formula has constant coefficient
$\sqrt{|\det g_0|}$. Substituting into the explicit formulas for connection,
curvature, star and divergence proves every assertion. Integration commutes
by its definition. At a halo point $a+h$, positive powers of $h$ disappear
under standard part, so these reductions take place at the real base point
$a$.
\end{proof}

\begin{proposition}[A quantitative support threshold]\label[proposition]{prop:threshold}
If $g=g_0+h$ and $\supp(h)\subseteq[\eta,\infty)$ for some $\eta>0$,
then the differences from the corresponding real quantities for $g^{-1}$,
$\Gamma$, $R$, $\Ric$, $\Scal$ and the volume density all have support in
$[\eta,\infty)$, locally in real coordinates.
\end{proposition}
\begin{proof}
Differentiation does not lower coefficient exponents. Multiplication by real
smooth coefficients does not change them. In every difference expansion,
including the inverse and binomial expansions, each term contains at least
one factor of $h$ or one of its real derivatives. It therefore has exponent
at least $\eta$. The expansions are strong by the positive-support lemma.
\end{proof}

The hypothesis is a support bound on a field and its chosen smooth
coefficient representation, not merely smallness of its values at a few
points. It also uses a nondegenerate real leading metric. If
$g=\operatorname{diag}(t^{2\eta},1,\ldots,1)$, the metric is perfectly
nondegenerate over $K$, but its standard part is degenerate and its inverse
contains $t^{-2\eta}$. The theorem intentionally excludes this situation.

\section{Examples at finite, infinitesimal, and infinite scales}\label{sec:examples}

\subsection{Anisotropic flat geometry}

For constant $\lambda_i\in\Gamma$ let
\[
 g=\sum_{i=1}^n t^{2\lambda_i}(\dd x^i)^2.
\]
It has
\[
 g^{-1}=\operatorname{diag}(t^{-2\lambda_i}),\quad
 \mu_g=t^{\lambda_1+\cdots+\lambda_n},\quad
 \Gamma=0,\quad R=0,
\]
and
\[
 \grad_g f=\sum_i t^{-2\lambda_i}D_i f\,D_i,
 \qquad \Delta_g f=\sum_i t^{-2\lambda_i}D_i^2f.
\]
The coordinate change $y^i=t^{\lambda_i}x^i$ makes the metric Euclidean
over $K$. It is generally not an ordinary real coordinate change. Different
component scales can be an exact change of units or frame, without creating
curvature.

\subsection{A smooth intrinsic metric with infinite curvature}\label{subsec:warped}

For $n\ge2$ and $\eps=t^\eta>0$ consider, on genuine $K$ coordinates,
\begin{equation}\label{eq:warped}
 g_\eps=(\dd x)^2+(x^2+\eps^2)(\dd y)^2
          +\sum_{j=3}^n(\dd z^j)^2.
\end{equation}
It is positive definite at every point. Its coefficients and inverse are
intrinsically smooth rational functions. The only nonzero Christoffel
coefficients in the two-dimensional factor are
\[
 \Gamma^x{}_{yy}=-x,\qquad
 \Gamma^y{}_{xy}=\Gamma^y{}_{yx}=\frac{x}{x^2+\eps^2}.
\]
Our curvature convention gives
\begin{equation}\label{eq:warpedcurv}
 \mathcal K_\eps=-\frac{\eps^2}{(x^2+\eps^2)^2},\qquad
 \Scal(g_\eps)=-\frac{2\eps^2}{(x^2+\eps^2)^2}.
\end{equation}
The extra flat factors do not change the scalar curvature. These formulas
follow directly from \eqref{eq:curvature}; the supplied symbolic program
checks the calculation independently.

At $x=0$, the sectional curvature is $-\eps^{-2}$, an infinite surreal.
At every fixed nonzero real $x$, it has leading term $-\eps^2/x^4$.
The inverse coefficient $(x^2+\eps^2)^{-1}$ has no smooth Hahn expansion
on a real neighborhood crossing zero. Thus this is an intrinsic metric, not
an admissible metric over that whole real base in the algebra
\eqref{eq:A}. There is no contradiction between those conclusions.

The inner coordinate $x=\eps X$ gives the two-dimensional factor
\[
 \eps^2\bigl((\dd X)^2+(1+X^2)(\dd y)^2\bigr),
 \qquad \mathcal K_\eps=-\eps^{-2}(1+X^2)^{-2}.
\]
A scale chart resolves the transition and exposes its curvature scale.
Replacing a small real parameter by a nonzero infinitesimal removes a zero
of this particular metric coefficient, but does not make the curvature
finite. Nondegeneracy, regular standard part and bounded invariants are
three different requirements.

\subsection{Radial fields and scale-independent flux}

In Euclidean $K^n$ let $r=(\sum_i(x^i)^2)^{1/2}$ and, away from $r=0$, set
\[
 X^i=c\,x^i r^{-n},\qquad c\in K.
\]
The local algebraic derivative of the positive square root gives
$D_i r=x^i/r$. Consequently
\[
 \Div X=c\left(nr^{-n}-nr^{-n-2}\sum_i(x^i)^2\right)=0.
\]
For $n\ne2$ this is the gradient of $c\,r^{2-n}/(2-n)$.
For $n=2$ a local logarithmic primitive may be chosen on an admitted positive
radius chart; no unrestricted global trigonometric or logarithmic function
is needed for the divergence calculation.

On the sphere parametrized by $x=R\theta$, $R\in K_{>0}$ and
$\theta\in S^{n-1}(\R)$, the pulled-back flux form is exactly
$c$ times the real unit-sphere area form: the factor $R^{1-n}$ in the field
cancels $R^{n-1}$ in the area. Define this parametric integral by its smooth
Hahn coefficients. Its value is
\[
 c\,\operatorname{area}(S^{n-1}),
\]
for infinitesimal, finite or infinite $R$. This is an explicit admissible
parametric calculation, not a claim that $\theta\mapsto R\theta$ is a
nonconstant fine-continuous map from a usual real sphere into full $\No^n$.
The central singularity at $r=0$ has not been assigned a value.

\section{Flows and parallel transport by positive-support recursion}\label{sec:flows}

Let $X_0$ be a real smooth vector field on an open $U\subset\R^n$. Suppose
its real flow $\phi_0(s,a)$ is defined smoothly for $(s,a)$ in a fixed real
flow domain containing $s=0$, and stays in $U$. Consider
\[
 X=X_0+\sum_{\gamma\in S}t^\gamma X_\gamma,
 \qquad S\subset\Gamma_{>0}\text{ well ordered},
\]
with all coefficient fields smooth on $U$.

\begin{theorem}[Near-real Hahn flows]\label{thm:flow}
There is a solution of
\[
 \partial_s\phi=(\ev X)(\phi),\qquad \phi(0,a)=a,
\]
of the form
\[
 \phi(s,a)=\phi_0(s,a)+\sum_{\gamma\in\langle S\rangle\setminus\{0\}}
 t^\gamma u_\gamma(s,a).
\]
Each coefficient is real smooth on the same real flow domain, the support is
well ordered, and the equation holds exactly in the Taylor--Hahn calculus.
The solution is unique among positive-support Hahn corrections with the
specified initial condition. The realized flow maps satisfy the local flow
composition law wherever all terms are defined.
\end{theorem}
\begin{proof}
Expand each coefficient vector field by its Taylor jet at $\phi_0$.
At exponent $\gamma$, the unknown $u_\gamma$ occurs only in the linear term
from $X_0$. Every other occurrence contains a strictly positive lower
exponent, either from $X-X_0$ or another correction. Thus
\[
 \partial_s u_\gamma
 =DX_0(\phi_0)u_\gamma+F_\gamma,
 \qquad u_\gamma(0,a)=0,
\]
where $F_\gamma$ is a finite expression in already constructed coefficients
and real derivatives. Finiteness even across all Taylor orders is
\cref{lem:neumann}. The ordinary inhomogeneous linear variational equation
has a unique smooth solution on the given real flow domain, by its
fundamental-matrix formula. This defines the coefficients by transfinite
recursion over $\langle S\rangle$.

The family is Hahn-supported by construction. Substitution back into the
Taylor--Hahn equation is strong and yields equality at every exponent.
For uniqueness against a second solution, take the union of the two positive
supports and the forcing support and apply the same least-differing-exponent
linear equation. The coefficient at that exponent must agree, a contradiction.
Finally, composing two flow maps is admissible by \eqref{eq:nonrealpull};
the usual autonomous equation and matching initial data, followed by this
uniqueness, prove the local flow law. The identities can be prolonged in the
initial coordinates by \cref{thm:prolongation}.
\end{proof}

The real flow domain is deliberately part of the hypothesis. This result
asserts neither completeness of arbitrary real fields nor existence of an
all-real-time flow for a field with infinite leading coefficients.
For example $X=t^{-\eta}D_1$ has the exact intrinsic solution
$x^1(s)=a^1+s t^{-\eta}$, which leaves the finite halo for nonzero real $s$.
The rescaled time $s=t^\eta\tau$ makes this motion finite.

For $n\ge2$, the field
\[
 X=t^\eta(x^1D_1-x^2D_2)
\]
has the exact flow
\[
 x^1(s)=a^1\sum_{m\ge0}\frac{(s t^\eta)^m}{m!},\qquad
 x^2(s)=a^2\sum_{m\ge0}\frac{(-s t^\eta)^m}{m!}.
\]
Its Jacobian determinant is $1$, by the strong exponential addition identity;
the other coordinates are fixed. This illustrates an exact infinitesimal
strain with exact volume conservation.

The same recursion proves parallel transport for a connection that is a
positive-support perturbation of a real connection along a near-real curve.
Geodesics reduce to a first-order system on the position--velocity space,
so \cref{thm:flow} applies near any real geodesic on its specified real
existence interval. Metric conservation then holds coefficientwise and on
the intrinsic realizations.

\section{Nonlinear field equations and a support-controlled lifting theorem}\label{sec:PDE}

The preceding ODE construction generalizes to a large class of equations for
scalar, vector or tensor fields. The point is not to assume that every real
PDE is solvable. It is to identify the exact real solvability input and then
prove the Hahn extension.

Let $E,F$ be real vector spaces of smooth sections with specified boundary,
normalization or gauge conditions. Suppose a finite-degree polynomial
differential operator has an expansion at a real field $u_0$,
\begin{equation}\label{eq:nonlinear}
 \mathcal P(u_0+w)=\mathcal P(u_0)+Lw+
       \sum_{r=2}^d B_r(w,\ldots,w),
\end{equation}
where $L:E\to F$ is real linear and each $B_r:E^r\to F$ is real multilinear.
The finite differential order is absorbed into these maps.

\begin{theorem}[Regular nonlinear Hahn lifting]\label{thm:lifting}
Assume $L:E\to F$ is bijective, with a fixed inverse taking the prescribed
real data to smooth sections with the prescribed conditions. Given a Hahn
forcing
\[
 f=\mathcal P(u_0)+\sum_{\gamma\in S}f_\gamma t^\gamma,
 \qquad S\subset\Gamma_{>0}\text{ well ordered},
\]
there is a unique solution $u=u_0+w$, with $w$ of positive Hahn support,
to the coefficientwise equation $\mathcal P(u)=f$. The correction $w$ has
support contained in $\langle S\rangle\setminus\{0\}$.
\end{theorem}
\begin{proof}
At $\gamma\in\langle S\rangle\setminus\{0\}$ define
\begin{equation}\label{eq:recursion}
 w_\gamma=L^{-1}\left(f_\gamma-
 \sum_{r=2}^d\ \sum_{\gamma_1+\cdots+\gamma_r=\gamma}
 B_r(w_{\gamma_1},\ldots,w_{\gamma_r})\right),
\end{equation}
where all $\gamma_i>0$ and missing forcing coefficients are zero. Every
$\gamma_i<\gamma$, so the recursion is well founded. Each inner sum is
finite by the support lemma. Since $L^{-1}$ maps real coefficient data to
real coefficient solutions, it cannot introduce Hahn exponents.
This constructs a Hahn family on the stated support; substitution and
finite coefficient comparison prove the equation.

For uniqueness, the union of two proposed positive supports is well ordered.
At the least exponent where the solutions differ, expand each nonlinear
difference as a telescoping sum of multilinear terms. Every such term
contains a difference of coefficients at a strictly smaller exponent and
therefore vanishes. The difference equation is $Lz=0$. Bijectivity gives
$z=0$, a contradiction.
\end{proof}

If $L$ has only a specified right inverse, the same recursion constructs a
solution with each correction in that right inverse's image. Uniqueness then
holds only with that normalization. A nontrivial cokernel creates compatibility
conditions. Gauge kernels and conservation identities in geometric PDEs
must be handled before applying the theorem.

\subsection{An exact nonlinear example in dimension \texorpdfstring{$n$}{n}}

On the ordinary flat torus $\mathbb T^n=(\R/2\pi\Z)^n$, take
\[
 \mathcal P(u)=(1-\Delta)u+u^2,
 \qquad f=\eps\cos x^1,\qquad \eps=t^\eta.
\]
The real operator $L=1-\Delta$ is bijective on smooth periodic functions:
its Fourier multiplier on mode $k\in\Z^n$ is $1+|k|^2$, and division by
it preserves rapid Fourier decay. The theorem gives an exact Hahn solution
$u=\sum_{m\ge1}u_m\eps^m$, with
\[
 u_1=L^{-1}\cos x^1,\qquad
 u_m=-L^{-1}\sum_{a+b=m}u_au_b\quad(m\ge2).
\]
The first terms are
\begin{align}\label{eq:examplePDE}
 u={}&\frac{\eps}{2}\cos x^1
 -\eps^2\left(\frac18+\frac1{40}\cos 2x^1\right)\notag\\
 &+\eps^3\left(\frac{11}{160}\cos x^1+
                   \frac1{800}\cos 3x^1\right)
 +\sum_{m\ge4}u_m\eps^m.
\end{align}
The supplied program computes six coefficients and verifies the residual
through that order with exact rational Fourier arithmetic. The all-orders
existence and uniqueness are furnished by the theorem, not by that finite
check. The solution is exact even in value groups where the sequence
$m\eta$ is not cofinal and ordinary partial-sum convergence fails.

\subsection{Tensor and constrained systems}

For tensor equations the real spaces $E,F$ may be smooth tensor sections and
$B_r$ can include contractions and covariant derivatives of a fixed real
background. For equations containing $g^{-1}$, the positive-support inverse
expansion gives an infinite Taylor nonlinearity rather than a polynomial.
The same coefficient recursion works when those multilinear terms have the
Neumann finite-word certificate. This additional certificate must be checked,
as must the invertibility of the gauge-fixed real linearization. The theorem
does not, by itself, establish global Einstein evolution, Navier--Stokes
regularity, or a singularity-removal result.

\section{Variational calculus and gauge-valued fields}\label{sec:variation}

\subsection{A scalar action and its Euler--Lagrange equation}

Let a real compact domain have an admissible metric and density, and let
$V$ be a polynomial potential with coefficients in $K$. For an admissible
scalar field set
\[
 \mathcal S[\phi]=\int^H
 \left(\tfrac12g^{ij}D_i\phi D_j\phi+V(\phi)\right)\vol_g.
\]
For a compactly supported variation $\psi$, take the coefficient of the
ordinary formal variation parameter $s$ in $\mathcal S[\phi+s\psi]$.
Leibniz's rule and Hahn Stokes give
\[
 \delta\mathcal S[\phi](\psi)
 =\int^H\bigl(-\Delta_g\phi+V'(\phi)\bigr)\psi\vol_g.
\]
The fundamental lemma remains valid in this setting: if this integral is zero
for every real smooth compactly supported $\psi$, then each coefficient of
the admissible residual density has zero integral against every real test
function, so each is zero by the real fundamental lemma. Hence stationarity
is equivalent to
\[
 -\Delta_g\phi+V'(\phi)=0.
\]
This is a precise formal-coefficient variational problem; its stationary
points and their existence are separate questions.

Define the stress tensor
\[
 T_{ij}=D_i\phi D_j\phi-
 g_{ij}\left(\tfrac12g^{ab}D_a\phi D_b\phi+V(\phi)\right).
\]
Metric compatibility, Hessian symmetry and the product rule yield
\[
 \nabla^iT_{ij}=(\Delta_g\phi-V'(\phi))D_j\phi.
\]
Thus the field equation implies covariant stress conservation exactly,
including all infinitesimal orders. This identity makes no positivity claim
when the metric has indefinite signature.

\subsection{Finite-rank gauge fields}

Let $A$ be a matrix-valued admissible one-form acting on a finite-rank
$K$-bundle. Its curvature and covariant exterior derivative are
\[
 F_A=dA+A\wedge A,\qquad
 d_A\beta=d\beta+A\wedge\beta-(-1)^{\deg\beta}\beta\wedge A.
\]
Direct expansion gives $d_AF_A=0$ and $d_A^2\beta=[F_A,\beta]$ in the
graded sense. Under an admissible frame transformation $P$ with admissible
inverse,
\[
 A'=P^{-1}AP+P^{-1}dP,\qquad F_{A'}=P^{-1}F_AP.
\]
The derivative of $P^{-1}P=I$ proves the transformation law. Positive-support
transformations $P=P_0(I+H)$ with real invertible $P_0$ are a concrete
admissible class.

In any dimension, the Yang--Mills-type equation $d_A(*F_A)=0$ is therefore
well defined whenever the metric star and products are admissible. As with
the scalar action, finite-rank algebra and variational identities do not
supply an automatic global existence theorem. Surcomplex coefficients can
be introduced by $K[i]$, with conjugation and Hermitian pairings specified
separately; they do not alter the real-coordinate support arguments.

\section{The independent scalar derivation}\label{sec:BM}

Spatial differentiation in this article fixes every $K$ constant. In
particular,
\[
 D_i\omega=0.
\]
The Berarducci--Mantova derivation $\partial$ on surreal scalars instead has
$\partial\omega=1$ and is strongly additive on summable normal forms
\cite{BM}. It describes differential-field structure, not the variation of
a constant scalar field across space. These two derivatives must not be
substituted for one another. The infinite-argument function and composition
calculus of omega-series \cite{BMGerms,Mantova26} is also additional structure,
not a universal replacement for coordinate calculus.

There is a compatible mixed construction when it is useful. For a given
Hahn field with smooth real coefficient functions, define
\[
 \partial F=\sum_\gamma f_\gamma\,\partial(t^\gamma).
\]
Strong additivity of the scalar derivation provides a well-ordered support
union and finite contributions. Coefficients are finite real linear
combinations of the $f_\gamma$, so they are smooth. The image may require a
larger workspace containing the derivative supports. On such admissible
expressions,
\[
 [D_i,\partial]F=0,
\]
because the $D_i$ act only on the real coefficient functions, while
$\partial$ acts only on the monomials.

An additional connection direction can be written
$\nabla_\partial=\partial+B$, with the semilinear Leibniz rule
$\nabla_\partial(aY)=\partial(a)Y+a\nabla_\partial Y$.
Together with $\nabla_i=D_i+\Gamma_i$ its mixed curvature is
\[
 [\nabla_\partial,\nabla_i]
 =\partial\Gamma_i-D_iB+[B,\Gamma_i].
\]
This is a differential-module extension of the theory. Interpreting it as an
extra physical time or scale direction would be a further modeling choice;
it is not forced by the definition of a surreal-valued tensor field.

\section{Computation, proof boundaries, and implementation}\label{sec:implementation}

\subsection{Representations and exactness}

A practical implementation should keep four objects separate: finite tensors
over an exact scalar field; sparse finite Hahn truncations; support-certified
infinite Hahn expressions; and coefficient operators on smooth or formal
real functions. A truncation is not an exact infinite series unless its tail
is separately represented or certified.

For a grid generated by positive exponents $\eta_1,\ldots,\eta_m$, a useful
normal form records a finite exponent vector, a real symbolic coefficient,
and a support certificate for allowed sums. Negative initial shifts are
stored separately. At arbitrary ordered-group rank, a cutoff $\gamma<B$
need not leave finitely many terms: for $0<\eta\ll B$, all $m\eta$ lie
below $B$. Algorithms must therefore use an explicit finite set of target
coefficients, a degree bound, or a proved left-finiteness condition. This
restriction is consistent with the repository's distinction between exact
denotation and effective support access \cite{Repo}.

\subsection{A minimal implementation architecture}

A suitable dependency order is: finite tensor algebra over an abstract real
closed field; Hahn support and strong evaluation; smooth coefficient
operators; the prolongation bridge; admissible inverses and metrics; tensor
calculus and curvature; coefficientwise chain integration; and finally
support-controlled equation solvers. A theorem about a metric inverse should
consume a unit certificate. A theorem about a substitution should consume a
joint support certificate. These are mathematical data, not comments that a
simplifier may silently discard.

For Lean, the finite part can be formulated over an abstract ordered field,
with additional square-root or real-closed hypotheses where needed. The Hahn
part should be proved separately and then instantiated through a verified
normal-form bridge. A global class of all surreals should not be forced into
a smaller universe by asserting that all its cuts or all its coefficients
form a set. Our proofs describe this layering; no new Lean implementation or
kernel verification is claimed here.

\subsection{Included finite verification program}

The archive includes \path{code/verify.py}. It independently checks the
warped-metric curvature, first-order metric inverse, the Lie-bracket Jacobi
identity for sample polynomial fields, $d^2=0$ for polynomial forms,
Hodge-star-square signs over dimensions $1$ through $6$ and every diagonal
signature, a radial homotopy identity, and six coefficients of the nonlinear
periodic example. All seven check groups passed, with 5,533 exact assertions.
The tested environment and outputs are recorded in \path{data/verification.txt}.

These checks are exact finite arithmetic and symbolic identities. They do
not prove real closedness, the support lemma, the all-orders recursion,
global cohomology, or the derivative bridge. Those are proved or explicitly
imported in the text. Conversely, successful typesetting of this article
has no bearing on mathematical correctness.

\section{What the theory establishes, and what it leaves as choices}\label{sec:conclusion}

The resulting theory has three genuine strengths. It accommodates every
finite tensor type in every finite dimension with exact infinite and
infinitesimal scalar values. It provides a coordinate-compatible calculus
with certified infinite operations, including actual intrinsic derivatives
on a large class of surreal-coordinate domains. And it converts specified
real linear solvability into exact nonlinear multiscale solutions without
requiring ordinary convergence of perturbation series.

Its boundaries are equally concrete. A bare map into $\No^n$ is not a
smooth field. Pointwise nondegeneracy is not a certificate that an inverse
belongs to a particular function algebra. The intrinsic topology of a fixed
workspace is not the topology inherited from the full proper class.
Coefficientwise integration is not an unrestricted measure theory on all
surreal subsets. Standard part preserves regular geometry but can destroy
rank in a multiscale metric. A field derivation on scalars is not a spatial
partial derivative. None of these distinctions disappears when $n$ is
specialized to three or four.

Several further developments can now be stated as precise mathematical
projects rather than undefined analogies: characterize larger intrinsic
function classes admitting coordinate gluing and integration; prove Hodge
and elliptic results for genuinely surreal leading metrics; classify scale
charts around rank-changing real limits; establish support certificates for
specific gauge-fixed nonlinear systems; and formalize the prolongation and
unit-admissibility interfaces. The constructions already proved here give a
working foundation for those projects, while retaining a sharp distinction
between exact scale representation and analytic or physical regularization.

\appendix
\section{Notation and hypothesis ledger}

\begin{longtable}{P{0.24\textwidth}P{0.68\textwidth}}
\toprule
Symbol or phrase & Meaning and scope\\
\midrule\endhead
$\No$ & The proper class of Conway surreal numbers; not a set-sized scalar type without a universe qualification.\\
$\Gamma$ & A nonzero set-sized divisible ordered additive subgroup of $\No$.\\
$t^\gamma$ & Conway monomial $\omega^{-\gamma}$, not an implicit use of surreal exponentiation.\\
$K$ & The real closed Hahn workspace $\R((t^\Gamma))$, with its intrinsic topology when derivatives are taken.\\
$\A_\Gamma(U)$ & Smooth real coefficient functions on one common real domain, with one well-ordered outer Hahn support.\\
$\A^{\mathrm{loc}}_\Gamma$ & The sheaf of local such presentations; globally uniform support is not assumed on noncompact bases.\\
$v$, $\vH$ & Minimum exponent of a scalar; minimum exponent of a nonzero coefficient function in a local field presentation.\\
$\halo U$ & $\st^{-1}(U)\subset K^n$; a nearstandard open set, not all of $K^n$.\\
$D_i$ & Coefficientwise spatial differentiation, fixing every constant in $K$.\\
$\partial_{y^i}$ & Intrinsic ordered-field derivative of a realized function; agrees with $\ev D_i$ only where the bridge applies.\\
$\partial$ & The optional Berarducci--Mantova scalar derivation, not a spatial derivative.\\
Admissible metric & Symmetric metric whose inverse lies in the chosen coefficient algebra; its density is also required when used.\\
Regular metric & $g=g_0+h$ with real nondegenerate $g_0$ and positive-support $h$.\\
$\Delta_g$, $\Delta_H$ & Divergence--gradient scalar Laplacian; $d\delta+\delta d$. On functions $\Delta_H=-\Delta_g$.\\
$\int^H$, $\int^{\mathrm{alg}}$ & Coefficientwise integration over real chains; polynomial moment integration on affine $K$-simplices.\\
Strong sum & Well-ordered support union and finitely many nonzero contributions at every exponent; no ordinary convergence assertion.\\
\bottomrule
\end{longtable}

\section{Repository and source audit}\label{sec:audit}

The repository was inspected read-only at commit
\begin{center}\small\texttt{\repoSHA}.\end{center}
The targeted review used the root \path{README.md}, \path{docs/README.md},
the opening scope of
\path{docs/foundations-and-computation/foundations/article.tex}, and the
normal-form, workspace, function-class, derivative and support conventions in
\path{docs/surcomplex/analysis/article.tex}. The documentation catalogue was
used to identify the nearby contours, differential-equations, spectral and
physics reports. This was not a line-by-line audit of every manuscript or a
local rebuild of the Lean project. The catalogue's report count is not used
as a claim that the entire repository, including newly added material, was
exhaustively surveyed.

The points carried forward are the use of increasing Hahn exponents,
set-sized divisible workspaces, common domains, strong summability, the
separation of scalar and coordinate derivations, and explicit statements of
formalization scope. The present article's tensor and PDE developments are
proved independently above. The repository itself describes its manuscripts
as research drafts with theorem-specific, rather than blanket, Lean coverage.

The external literature check included the original Berarducci--Mantova
paper, the general-base-field calculus of Bertram--Gl\"ockner--Neeb, the
van den Dries--Ehrlich paper, Poonen's discussion of Hahn units, and classical
smooth and de Rham geometry sources. Two sources relevant to the cutoff date
are Poonen's January 2026 revision on units and Mantova's January 2026
Taylor-approximation preprint. Their roles are identified in the bibliography;
neither is represented as a theorem about all tensor fields on all of
$\No^n$.

\section{A proof-dependency summary}

The support-sensitive claims have short, explicit dependency chains.
The field and size setup uses normal forms and the classical Hahn-field
closure theorem. Differential algebra uses the finite-pair support lemma.
Prolongation adds positive-support substitution and a valuation remainder
estimate. Metrics add only a certified geometric inverse and a certified
square root. Curvature identities are finite differential-algebra identities.
Stokes and Poincar\'e use fixed real coefficient operators. Compact de Rham
base change uses finite real cohomology and support-preserving linear
splittings; the Hodge refinement imports the real Hodge theorem. Nonlinear
lifting uses a specified real linear inverse plus positive-support recursion.
No circular reliance on the desired surreal existence theorem occurs in
these constructions.

\begin{thebibliography}{99}
\small
\bibitem{Repo}
Vladimir Reshetnikov, \emph{Surreal}, repository documentation, inspected at
commit \texttt{\repoSHA}, September 22, 2026.
\url{https://github.com/VladimirReshetnikov/Surreal}.
The root README and \path{docs/README.md} give the implementation and
manuscript scope; these are project sources, not a blanket verification claim.

\bibitem{RepoFound}
\emph{Foundations for Surreal and Surcomplex Mathematics}, Surreal project,
\path{docs/foundations-and-computation/foundations/article.tex}, same pinned
revision as \cite{Repo}. Targeted use: sizes, workspaces and the fine-topology
warning.

\bibitem{RepoAnalysis}
\emph{Surcomplex Analysis}, Surreal project,
\path{docs/surcomplex/analysis/article.tex}, same pinned revision as
\cite{Repo}. Targeted use: monomials, coherent coefficient fields,
summability and the distinction between derivative operators.

\bibitem{Gonshor}
Harry Gonshor, \emph{An Introduction to the Theory of Surreal Numbers},
London Mathematical Society Lecture Note Series 110, Cambridge University
Press, 1986.
\url{https://www.cambridge.org/core/books/an-introduction-to-the-theory-of-surreal-numbers/312AE504A3E88E804054BFB390446374}.
Background for normal forms and surreal algebra.

\bibitem{vdDE}
Lou van den Dries and Philip Ehrlich, ``Fields of surreal numbers and
exponentiation,'' \emph{Fundamenta Mathematicae} 167 (2001), 173--188.
DOI: \href{https://doi.org/10.4064/fm167-2-3}{10.4064/fm167-2-3}.
See also their published erratum, \emph{Fundamenta Mathematicae} 168 (2001).
Used as background, not as an unrestricted transfer theorem for our calculus.

\bibitem{Poonen93}
Bjorn Poonen, ``Maximally complete fields,'' \emph{L'Enseignement
Math\'ematique} (2) 39 (1993), 87--106.
\url{https://math.mit.edu/~poonen/papers/amsval.pdf}.
Used for the algebraic-closure theorem for Hahn fields with algebraically
closed residue field and divisible value group.

\bibitem{Neumann}
B. H. Neumann, ``On ordered division rings,'' \emph{Transactions of the
American Mathematical Society} 66 (1949), 202--252.
Classical source for the well-ordered support and positive-word finiteness
arguments.

\bibitem{Poonen26}
Bjorn Poonen, ``Units in Hahn--Mal'cev--Neumann rings,'' author manuscript,
January 14, 2026.
\url{https://math.mit.edu/~poonen/papers/malcev.pdf}.
Theorems 1.1 and 4.1 give the positive-support unit theorem and its explicit
geometric-series form.

\bibitem{BGN}
Wolfgang Bertram, Helge Gl\"ockner and Karl-Hermann Neeb,
``Differential calculus, manifolds and Lie groups over arbitrary infinite
fields,'' preprint, 2003, arXiv:math/0303300.
\url{https://arxiv.org/abs/math/0303300}.
A general framework showing why calculus over a non-real field requires an
explicit differentiability theory.

\bibitem{Lee}
John M. Lee, \emph{Introduction to Smooth Manifolds}, second edition,
Graduate Texts in Mathematics 218, Springer, 2013.
DOI: \href{https://doi.org/10.1007/978-1-4419-9982-5}{10.1007/978-1-4419-9982-5}.
Background for real smooth jets, tensor bundles, forms and Stokes.

\bibitem{Petersen}
Peter Petersen, \emph{Manifolds, Transversality, and de Rham Cohomology},
author lecture notes, accessed September 22, 2026, especially Chapters 6--7.
\url{https://www.math.ucla.edu/~petersen/manifolds.pdf}.
Background for real tensor calculus and de Rham arguments.

\bibitem{Warner}
Frank W. Warner, \emph{Foundations of Differentiable Manifolds and Lie
Groups}, Graduate Texts in Mathematics 94, Springer, 1983, Chapter 6.
DOI: \href{https://doi.org/10.1007/978-1-4757-1799-0}{10.1007/978-1-4757-1799-0}.
The real compact-manifold Hodge theorem is the imported analytic input to
\cref{cor:hodge}.

\bibitem{BM}
Alessandro Berarducci and Vincenzo Mantova, ``Surreal numbers, derivations
and transseries,'' \emph{Journal of the European Mathematical Society} 20
(2018), 339--390. DOI: \href{https://doi.org/10.4171/JEMS/769}{10.4171/JEMS/769}.
\url{https://arxiv.org/abs/1503.00315}.
Used only for the optional scalar-derivation direction.

\bibitem{BMGerms}
Alessandro Berarducci and Vincenzo Mantova, ``Transseries as germs of surreal
functions,'' arXiv:1703.01995.
\url{https://arxiv.org/abs/1703.01995}.
A separate infinite-argument composition calculus, not a claim about arbitrary
surreal-coordinate tensor fields.

\bibitem{Mantova26}
Vincenzo Mantova, ``Monotonicity and a Taylor approximation theorem for
transseries,'' preprint, January 12, 2026, arXiv:2601.07747.
\url{https://arxiv.org/abs/2601.07747}.
Recent context for omega-series composition and its Taylor theory; not used
in the proof of the smooth finite-halo prolongation theorem.
\end{thebibliography}
\end{document}
